%%%%%%%%%%%%%%%%%%%%%%%%%%%%%%%%%%%%%%%%%%%%%%%%%%%%%%%%%%%%%%
% BEGIN RESEARCH
\begin{rSection}{Professional EXPERIENCE}

    \setlength{\parskip}{1pt}
    \setlength{\baselineskip}{3pt}

\textbf{Data Engineering Intern, HAOCHENG TECH HOLDING} \hfill July 2022 -- August 2022 \\
\textit{Shenzhen, Guangdong}
\begin{itemize}
\vspace{-5pt}
    \item Architected and deployed \textbf{PostgreSQL} database schema using \textbf{DataGrip}, implementing B-tree and composite indexes on video metadata fields. Achieved a \textbf{15x performance improvement} by reducing search operation latency from 30 minutes to 2 minutes.
    
    \item Engineered end-to-end data pipeline and developed classification models integrating classical \textbf{machine learning} algorithms (\textbf{Random Forest} and \textbf{Logistic Regression}) to predict material classification on a large CSV dataset, achieving 87.3\% accuracy.

        \item Applied \textbf{NLP} techniques including \textbf{TF-IDF} vectorization and \textbf{text preprocessing} (tokenization) using \textbf{NLTK} and \textbf{scikit-learn} to classify and tag \textbf{excel} informaton for downstream search optimization.


\end{itemize}

\textbf{LLM Agent Implementation and Fine-Tuning in Data Science}\hfill February 2025 - May 2025
\begin{itemize}
\vspace{-5pt}
\item Developed \textbf{prompt engineering pipeline} implementing data cleaning, normalization, and serialization to \textbf{JSON} format with accompanying visualization generation, building a 500+ example prompt corpus with interpolation curve ground truth labels for \textbf{supervised few-shot learning} and model benchmarking.
    \item Implemented LLM agent onto \textbf{SSH}  to train the classification \textbf{vision-Transformer }and fine-tuning by \textbf{LoRA}, achieved
    90.4\% accuracy for astronomical data.

\item Generated and trained \textbf{Deep Learning }\textbf{LeNet-5 neural network} on MNIST dataset (60K images) achieving 80.3\% test accuracy to establish a benchmark classifier.
\end{itemize}

    
\textbf{Remote GPU Server Implementation and Numerical Computation} \hfill July 2023 - September 2023

\vspace{-5pt}

    
\begin{itemize}

    \item Develop Python code to realize the FFT matrix acceleration algorithm
on a remote server with \textbf{GPU} using \textbf{pytorch}, \textbf{tensorflow} and \textbf{cupy}, achieved approximately 20\% performance improvement.
\end{itemize}





\end{rSection} 
%%%%%%%%%%%%%%%%%%%%%%%%%%%%%%%%%%%%%%%%%%%%%%%%%%%%%%%%%%%%%%
% END RESEARCH